\section{The rational-function example}

\begin{definition}[Transcendental unit-circle parameters]
\label{def:rational-circle-parameters}
\source{chen-moriwaki-adelic:page-0189}
\lean{ArakelovGeometry.AdelicCurve.RationalFunction.circlePoint, ArakelovGeometry.AdelicCurve.RationalFunction.circlePoint_eq_exp, ArakelovGeometry.AdelicCurve.RationalFunction.continuous_circlePoint, ArakelovGeometry.AdelicCurve.RationalFunction.archimedeanDomain, ArakelovGeometry.AdelicCurve.RationalFunction.ArchimedeanPlace, ArakelovGeometry.AdelicCurve.RationalFunction.archimedeanMeasure}
\leanok
Write $z(t)=\exp(2\pi it)$ and
$T=\{t\in[0,1]:z(t)\text{ is transcendental over }\mathbb Q\}$.
The map $z:\mathbb R\to\mathbb C$ is continuous. Equip $T$ with the
trace Borel $\sigma$-algebra and induced Lebesgue measure.
\end{definition}
\begin{proof}
Multiplication by $2\pi i$ and the complex exponential are continuous,
so their composition $z$ is continuous. The trace $\sigma$-algebra and
the induced measure are the subspace constructions for the inclusion
$T\hookrightarrow\mathbb R$.
\end{proof}

\begin{lemma}[Full measure of the transcendental parameters]
\label{lem:rational-transcendental-circle-conull}
\source{chen-moriwaki-adelic:page-0189}
\uses{def:rational-circle-parameters}
\lean{ArakelovGeometry.AdelicCurve.RationalFunction.countable_algebraic_circle_parameters, ArakelovGeometry.AdelicCurve.RationalFunction.measurableSet_archimedeanDomain, ArakelovGeometry.AdelicCurve.RationalFunction.volume_archimedeanDomain}
\leanok
The set $\{t\in\mathbb R:z(t)\text{ is algebraic over }\mathbb Q\}$
is countable. Consequently $[0,1]\setminus T$ is countable, $T$ is Borel
measurable, and its induced Lebesgue measure is a probability measure.
\end{lemma}
\begin{proof}
The algebraic complex numbers form a countable set: there are countably many
nonzero rational polynomials, and each has finitely many roots. The map
$t\mapsto z(t)$ is injective on $(0,1]$, because equality of its values implies
that the difference of its arguments is an integer, and this difference has
absolute value less than one. Thus its inverse image of the algebraic numbers
inside $(0,1]$ is countable. Adding the endpoint $0$ still gives a countable
set. By periodicity the entire exceptional set is contained in the union
of its integer translates, which is again countable.
Countable subsets of the real line are Borel and Lebesgue null. Removing
this set from the Borel interval $[0,1]$, of measure one, proves the claim.
\end{proof}

\begin{definition}[Archimedean rational-function places]
\label{def:rational-archimedean-place}
\source{chen-moriwaki-adelic:page-0189}
\uses{def:rational-circle-parameters}
\lean{ArakelovGeometry.AdelicCurve.RationalFunction.archimedeanEval, ArakelovGeometry.AdelicCurve.RationalFunction.archimedeanAbv, ArakelovGeometry.AdelicCurve.RationalFunction.archimedeanAbv_apply, ArakelovGeometry.AdelicCurve.RationalFunction.archimedeanAbv_rat, ArakelovGeometry.AdelicCurve.RationalFunction.not_isNonarchimedean_archimedeanAbv}
\leanok
For $t\in T$, evaluation at $z(t)$ defines a field homomorphism
$e_t:\mathbb Q(T)\longrightarrow\mathbb C$. Set $|f|_t=|e_t(f)|$.
This is an archimedean absolute value extending the usual absolute value of $\mathbb Q$.
If $p_f,q_f$ are the canonical numerator and denominator of $f$, then
$|f|_t=|p_f(z(t))/q_f(z(t))|$.
\end{definition}
\begin{proof}
Transcendence says that no nonzero rational polynomial vanishes at $z(t)$.
Polynomial evaluation is therefore injective and extends uniquely to the
fraction field by $e_t(p/q)=p(z(t))/q(z(t))$. This is a field homomorphism.
Pulling back the complex absolute value gives an absolute value, and
evaluation fixes rational constants, which proves the extension assertion.
It is archimedean because $|2|_t=2>1=\max\{|1|_t,|1|_t\}$.
\end{proof}

\begin{lemma}[Polynomial logarithms on the transcendental circle]
\label{lem:rational-polynomial-circle-log-integrable}
\source{chen-moriwaki-adelic:page-0189}
\uses{def:rational-circle-parameters,lem:rational-transcendental-circle-conull}
\lean{ArakelovGeometry.AdelicCurve.RationalFunction.integrable_log_polynomial_circle}
\leanok
For every rational polynomial $p$, the function
$t\mapsto\log|p(z(t))|$ is integrable on $T$, with the logarithm extended
by zero at zero.
\end{lemma}
\begin{proof}
The zero polynomial gives the zero function. For a nonzero polynomial,
the logarithm of the absolute value of a nonzero
complex polynomial is integrable on the unit circle. Indeed, factor the
polynomial as $a\prod_j(X-\alpha_j)$; outside its finitely many zeros,
its logarithm is $\log|a|+\sum_j\log|z-\alpha_j|$.
Changing its values at the finitely many zeros does not affect integrability.
Away from those zeros the logarithm is continuous and bounded on compact arcs. A root off the
circle contributes a continuous function. For a root
$\alpha_j=e^{i\theta_j}$ on the circle, write $u=\theta-\theta_j$ modulo
$2\pi$, with $|u|\leq\pi$. Then
\[
\frac{2}{\pi}|u|\leq|e^{iu}-1|=2|\sin(u/2)|\leq|u|.
\]
Consequently $\log|e^{iu}-1|-\log|u|$ is bounded for $u\ne0$,
while $\int_0^\varepsilon|\log u|\,du$ is finite, as follows from the
antiderivative $u\log u-u$. Thus each root contributes an integrable
function and their finite sum is integrable. The linear change of
variables $\theta=2\pi t$ gives integrability on $[0,1]$.
Restriction to the measurable subset $T$ preserves integrability.
\end{proof}

\begin{lemma}[Integrability on the archimedean component]
\label{lem:rational-archimedean-log-integrable}
\source{chen-moriwaki-adelic:page-0189}
\uses{def:rational-archimedean-place,lem:rational-polynomial-circle-log-integrable}
\lean{ArakelovGeometry.AdelicCurve.RationalFunction.measurable_log_archimedeanAbv, ArakelovGeometry.AdelicCurve.RationalFunction.integrable_log_archimedeanAbv}
\leanok
For every nonzero $f\in\mathbb Q(T)$, the function
$t\mapsto\log|f|_t$ is measurable and integrable on $T$.
\end{lemma}
\begin{proof}
Write $f=p/q$ with nonzero rational polynomials $p,q$. Each function
$t\mapsto\log|p(z(t))|$ is Borel measurable, since polynomial evaluation,
the exponential, the norm, and the real logarithm extended by zero at zero
are Borel measurable. The preceding lemma proves integrability of the
numerator and denominator logarithms. On $T$, neither $p(z(t))$ nor
$q(z(t))$ vanishes, so
$\log|f|_t=\log|p(z(t))|-\log|q(z(t))|$.
The difference is measurable and integrable.
\end{proof}

\begin{definition}[The archimedean component adelic structure]
\label{def:rational-archimedean-structure}
\source{chen-moriwaki-adelic:page-0189}
\uses{def:rational-archimedean-place,lem:rational-transcendental-circle-conull,lem:rational-archimedean-log-integrable}
\lean{ArakelovGeometry.AdelicCurve.RationalFunction.archimedeanAdelicStructure}
\leanok
The field $\mathbb Q(T)$, the probability space $T$ with induced Lebesgue
measure, and the absolute values $|\cdot|_t$ form an adelic structure.
\end{definition}
\begin{proof}
The component places are absolute values by construction, and the required
measurability and integrability of logarithms are precisely the preceding
lemma. The measure has mass one by the full-measure lemma.
\end{proof}

\subsection{The finite components over the rational function field}

Put $K=\mathbb Q(T)$ and let $\Omega_h$ be the closed points of
$\mathbb A^1_{\mathbb Q}$. We equip $\Omega_h$ with its discrete
$\sigma$-algebra and counting measure, as in
\cite[\S3.2.5]{chen-moriwaki}.

\begin{definition}[The horizontal absolute values]
\label{def:rational-horizontal-place}
\source{chen-moriwaki-adelic:page-0188}
\lean{ArakelovGeometry.AdelicCurve.RationalFunction.HorizontalPlace,ArakelovGeometry.AdelicCurve.RationalFunction.horizontalMeasurableSpace,ArakelovGeometry.AdelicCurve.RationalFunction.horizontalPolynomial,ArakelovGeometry.AdelicCurve.RationalFunction.horizontalPolynomial_isPrimitive,ArakelovGeometry.AdelicCurve.RationalFunction.horizontalPolynomial_ne_zero,ArakelovGeometry.AdelicCurve.RationalFunction.horizontalPolynomial_span,ArakelovGeometry.AdelicCurve.RationalFunction.horizontalHeight,ArakelovGeometry.AdelicCurve.RationalFunction.one_le_horizontalHeight,ArakelovGeometry.AdelicCurve.RationalFunction.horizontalAbv,ArakelovGeometry.AdelicCurve.RationalFunction.horizontalAbv_apply,ArakelovGeometry.AdelicCurve.RationalFunction.horizontalAbv_rat}
\leanok
For $x\in\Omega_h$, choose a generator of the corresponding nonzero prime
ideal of $\mathbb Q[T]$, clear its denominators, and take its primitive
part $F_x\in\mathbb Z[T]$. Then $F_x\ne0$ and generates the same ideal
over $\mathbb Q$. Equip $\Omega_h$ with the discrete $\sigma$-algebra.
Let $H(x)=M(F_x)$ be its Mahler measure. We have $H(x)\geq1$.
The horizontal absolute value is
\[
 |0|_x=0,\qquad |f|_x=H(x)^{-\operatorname{ord}_x(f)}\quad(f\ne0).
\]
It restricts to the trivial absolute value on $\mathbb Q$; when
$H(x)=1$, it is the trivial absolute value on $K$ itself.
\end{definition}
\begin{proof}
Clearing denominators multiplies a generator by a nonzero integer.
Taking the primitive part divides it by its nonzero integral content.
Both scalars are units in $\mathbb Q[T]$, so the resulting polynomial
still generates the ideal of $x$, is nonzero, and is primitive.
For every nonzero integral polynomial the Mahler measure is at least
the absolute value of its leading coefficient, hence at least one.
If $H(x)>1$, the valuation inequalities and multiplicativity of
$\operatorname{ord}_x$ imply the nonarchimedean triangle inequality and
multiplicativity of the displayed absolute value. If $H(x)=1$, the
displayed formula is the trivial absolute value and satisfies the same
axioms directly. Finally a nonzero constant is a unit in
$\mathbb Q[T]$, so its order at $x$ is zero and its absolute value is one.
\end{proof}

\begin{definition}[The prime Gauss absolute values]
\label{def:rational-prime-place}
\source{chen-moriwaki-adelic:page-0188,chen-moriwaki-adelic:page-0189}
\lean{ArakelovGeometry.AdelicCurve.RationalFunction.rationalPadicAbv,ArakelovGeometry.AdelicCurve.RationalFunction.primeMeasurableSpace,ArakelovGeometry.AdelicCurve.RationalFunction.isNonarchimedean_rationalPadicAbv,ArakelovGeometry.AdelicCurve.RationalFunction.primePolynomialValuation,ArakelovGeometry.AdelicCurve.RationalFunction.primePolynomialValuation_ne_zero,ArakelovGeometry.AdelicCurve.RationalFunction.primeValuation,ArakelovGeometry.AdelicCurve.RationalFunction.primeAbv,ArakelovGeometry.AdelicCurve.RationalFunction.primeAbv_polynomial,ArakelovGeometry.AdelicCurve.RationalFunction.isNonarchimedean_primeAbv,ArakelovGeometry.AdelicCurve.RationalFunction.primeAbv_num_div_denom,ArakelovGeometry.AdelicCurve.RationalFunction.primeAbv_rat}
\leanok
For a rational prime $p$ and a polynomial $P=\sum_i a_iT^i$, put
\[
 \lVert P\rVert_p=\max_i|a_i|_p,
\]
where the maximum is zero when $P=0$ and $|\cdot|_p$ is the usual
nonarchimedean rational $p$-adic absolute value. Equip the prime index
set with the discrete $\sigma$-algebra. The Gauss norm, regarded as a
valuation with values in $\mathbb R_{\geq0}$, is nonzero on every
nonzero polynomial. It defines a nonarchimedean absolute value on
$\mathbb Q[T]$ and extends to $K$ by
\[
 |P/Q|_p=\lVert P\rVert_p/\lVert Q\rVert_p\qquad(Q\ne0).
\]
Its restriction to the constant field is the original rational
$p$-adic absolute value.
\end{definition}
\begin{proof}
The coefficientwise nonarchimedean inequality gives
$\lVert P+Q\rVert_p\leq\max(\lVert P\rVert_p,\lVert Q\rVert_p)$;
nonnegativity and vanishing only on the zero polynomial are immediate.
For nonzero $P,Q$, choose the least indices $i,j$ attaining their
respective coefficient maxima. In the coefficient of degree $i+j$ of
$PQ$, the term $a_ib_j$ has absolute value
$\lVert P\rVert_p\lVert Q\rVert_p$. Every other term has strictly
smaller absolute value: a smaller first index gives a smaller first
factor, and a larger first index gives a second index smaller than $j$.
The strong triangle inequality then shows that this coefficient has
absolute value exactly $\lVert P\rVert_p\lVert Q\rVert_p$.
All coefficients are bounded above by that product, so the Gauss norm
is multiplicative. The zero cases follow immediately.
If $P/Q=P'/Q'$ with nonzero denominators, then $PQ'=P'Q$;
multiplicativity and division by the nonzero denominator norms prove
that the quotient formula is independent of the presentation.
Applying this formula after putting two fractions over a common
denominator proves the nonarchimedean inequality on $K$.
It also proves multiplicativity and vanishing precisely at zero.
For a constant polynomial the Gauss norm is its single coefficient's
absolute value, which proves the restriction assertion.
\end{proof}

\begin{lemma}[Integrability at the rational primes]
\label{lem:rational-prime-log-integrable}
\source{chen-moriwaki-adelic:page-0188,chen-moriwaki-adelic:page-0189}
\uses{def:rational-prime-place}
\lean{ArakelovGeometry.AdelicCurve.RationalFunction.rationalPadicAbv_eq_one_of_large,ArakelovGeometry.AdelicCurve.RationalFunction.finite_rationalPadicAbv_ne_one,ArakelovGeometry.AdelicCurve.RationalFunction.finite_primeAbv_polynomial_ne_one,ArakelovGeometry.AdelicCurve.RationalFunction.finite_support_log_primeAbv,ArakelovGeometry.AdelicCurve.RationalFunction.integrable_log_primeAbv,ArakelovGeometry.AdelicCurve.RationalFunction.primeAdelicStructure}
\leanok
For $a=m/n\in\mathbb Q^\times$ in lowest terms, $|a|_p=1$ whenever
$p>\max(|m|,n)$, so only finitely many primes have $|a|_p\ne1$.
For every nonzero polynomial $P\in\mathbb Q[T]$, only finitely many
primes have $\lVert P\rVert_p\ne1$.
For $f\in K^\times$, the function $p\mapsto\log|f|_p$ has finite
support on the set of rational primes. It is therefore integrable for
counting measure and defines the prime adelic structure on $K$.
\end{lemma}
\begin{proof}
For $a=m/n\in\mathbb Q^\times$ in lowest terms, a prime exceeding
$\max(|m|,n)$ divides neither $m$ nor $n$, so $|a|_p=1$.
There are only finitely many smaller primes. Thus each nonzero rational
coefficient has absolute value one outside a finite set of primes.
For a nonzero polynomial, take the union of these exceptional sets
over its finitely many nonzero coefficients. Outside this finite
union all nonzero coefficients have norm one, so their maximum is one.
Finally write $f=P/Q$ with $P,Q\ne0$. Outside the union of the finite
exceptional sets for $P$ and $Q$, its absolute value is $1/1=1$.
The logarithm vanishes there. On the discrete prime space measurability
is automatic, and counting measure integrates its absolute value as a
finite sum. This proves all the adelic integrability axioms.
\end{proof}

\begin{lemma}[Integrability at the horizontal places]
\label{lem:rational-horizontal-log-integrable}
\source{chen-moriwaki-adelic:page-0188,chen-moriwaki-adelic:page-0189}
\uses{def:rational-horizontal-place}
\lean{ArakelovGeometry.AdelicCurve.RationalFunction.finite_support_log_horizontalAbv,ArakelovGeometry.AdelicCurve.RationalFunction.integrable_log_horizontalAbv,ArakelovGeometry.AdelicCurve.RationalFunction.horizontalAdelicStructure}
\leanok
For $f\in K^\times$, the function $x\mapsto\log|f|_x$ has finite
support on $\Omega_h$. Consequently it is integrable for counting
measure and defines the horizontal adelic structure on $K$.
\end{lemma}
\begin{proof}
Write $f=P/Q$ with nonzero $P,Q\in\mathbb Q[T]$.
Outside the finite union of the prime divisors of $P$ and $Q$,
$\operatorname{ord}_x(f)=0$. The preceding formula then gives
$|f|_x=1$, so its logarithm is zero.
Every function on the discrete space $\Omega_h$ is measurable.
Since each singleton has counting measure one, the integral of the
absolute value of a finitely supported real function is a finite sum
of finite real numbers. This proves integrability, and hence every
axiom of the asserted adelic structure.
\end{proof}


\begin{definition}[The adelic structure on $\mathbb Q(T)$]
\label{def:rational-function-adelic-curve}
\source{chen-moriwaki-adelic:page-0188,chen-moriwaki-adelic:page-0189,chen-moriwaki-adelic:page-0220}
\uses{def:rational-horizontal-place,lem:rational-horizontal-log-integrable,def:rational-prime-place,lem:rational-prime-log-integrable,def:rational-archimedean-structure}
\lean{ArakelovGeometry.AdelicCurve.RationalFunction.finiteAdelicStructure, ArakelovGeometry.AdelicCurve.RationalFunction.adelicStructure, ArakelovGeometry.AdelicCurve.RationalFunction.adelicCurve, ArakelovGeometry.AdelicCurve.RationalFunction.adelicStructure_horizontal, ArakelovGeometry.AdelicCurve.RationalFunction.adelicStructure_prime, ArakelovGeometry.AdelicCurve.RationalFunction.adelicStructure_archimedean, ArakelovGeometry.AdelicCurve.RationalFunction.adelicStructure_measure}
\leanok
Let $\Omega_h$ be the closed points of $\mathbb A^1_{\mathbb Q}$ and
$\mathcal P$ the rational primes. The parameter space
\[
 \Omega=(\Omega_h\amalg\mathcal P)\amalg T
\]
has counting measure on its first two components and Lebesgue measure on
$T$. Equip $\mathbb Q(T)$ with the horizontal Mahler absolute values, the
$p$-adic Gauss absolute values, and the transcendental evaluation absolute
values on these three components respectively. This defines an adelic
structure $S$, obtained by first amalgamating the horizontal and prime
component structures and then amalgamating with the archimedean component.
\end{definition}
\begin{proof}
The component absolute values have already been constructed. Every function
on either discrete component is measurable. The preceding archimedean
integrability lemma gives measurability on the third component. For a nonzero
rational function the logarithmic evaluation has finite support on each
discrete component, hence is integrable there for counting measure. It is
integrable on the third component by the same lemma. Integrability and
measurability on a disjoint union are equivalent to the corresponding
properties on its components. These are exactly the adelic-structure axioms.
\end{proof}

\begin{remark}
The horizontal component is the one specified in Subsection~3.2.5, namely
the closed points of $\mathbb P^1_{\mathbb Q}\setminus\{\infty\}$.
The phrase ``closed points of $\mathbb P^1_{\mathbb Q}$'' in
Example~3.7.2(3) omits this exclusion. All absolute values and the parameter
measure here use the construction in Subsection~3.2.5.
\end{remark}

\begin{definition}[The augmented rational adelic curve]
\label{def:augmented-rational-adelic-curve}
\source{chen-moriwaki-adelic:page-0220}
\uses{def:rational-function-adelic-curve,def:rational-horizontal-place,def:rational-prime-place}
\lean{ArakelovGeometry.AdelicCurve.RationalFunction.rationalFiniteAdelicStructure, ArakelovGeometry.AdelicCurve.RationalFunction.rationalInfiniteAdelicStructure, ArakelovGeometry.AdelicCurve.RationalFunction.augmentedRationalAdelicStructure, ArakelovGeometry.AdelicCurve.RationalFunction.augmentedRationalAdelicCurve, ArakelovGeometry.AdelicCurve.RationalFunction.augmentedRationalAdelicStructure_horizontal, ArakelovGeometry.AdelicCurve.RationalFunction.augmentedRationalAdelicStructure_prime, ArakelovGeometry.AdelicCurve.RationalFunction.augmentedRationalAdelicStructure_infinite, ArakelovGeometry.AdelicCurve.RationalFunction.augmentedRationalAdelicStructure_measure}
\leanok
Let $\widetilde\Omega_{\mathbb Q}=(\Omega_h\amalg\mathcal P)\amalg\{\infty\}$,
with counting measure, and attach the trivial absolute value of $\mathbb Q$
to each $x\in\Omega_h$, the standard $p$-adic absolute value to each
$p\in\mathcal P$, and the usual real absolute value to $\infty$.
This defines an adelic curve $\widetilde S_{\mathbb Q}$.
Its finite component is the restriction of the finite component of $S$
along $\mathbb Q\hookrightarrow\mathbb Q(T)$; the full structure is its
amalgamation with the singleton infinite component.
It is the amalgamation of the standard rational adelic curve with
counting-measure copies of the trivial absolute value, with the three
components regrouped as displayed.
\end{definition}
\begin{proof}
The logarithmic functions are measurable since every component is discrete.
On $\Omega_h$ every nonzero rational has absolute value one, so its logarithm
vanishes. On the primes the logarithmic function is the restriction of the
corresponding Gauss logarithmic function for the constant rational function,
and thus is integrable. The infinite component is one point of mass one and
every real function there is integrable. This proves the adelic axioms.
The listed measures and absolute values coincide component by component
with the stated amalgamation; regrouping does not change them.
\end{proof}

\begin{example}[The rational-function morphism]
\label{ex:adelic-morphism-rational-function}
\source{chen-moriwaki-adelic:page-0220}
\uses{def:adelic-morphism,def:rational-function-adelic-curve,def:augmented-rational-adelic-curve,lem:rational-transcendental-circle-conull,lem:probability-collapse}
\lean{ArakelovGeometry.AdelicCurve.RationalFunction.morphismToRationals, ArakelovGeometry.AdelicCurve.RationalFunction.morphismToRationals_fieldMap, ArakelovGeometry.AdelicCurve.RationalFunction.morphismToRationals_baseMap_finite, ArakelovGeometry.AdelicCurve.RationalFunction.morphismToRationals_baseMap_archimedean, ArakelovGeometry.AdelicCurve.RationalFunction.fibreAverage, ArakelovGeometry.AdelicCurve.RationalFunction.fibreAverage_finite, ArakelovGeometry.AdelicCurve.RationalFunction.fibreAverage_infinite, ArakelovGeometry.AdelicCurve.RationalFunction.integrable_fibreAverage, ArakelovGeometry.AdelicCurve.RationalFunction.morphismToRationals_fibreIntegral_toL1}
\leanok
There is a morphism $\alpha:S\to\widetilde S_{\mathbb Q}$ whose field map
is the inclusion $\mathbb Q\hookrightarrow\mathbb Q(T)$, whose parameter
map fixes $\Omega_h\amalg\mathcal P$ and collapses $T$ to $\infty$, and
whose fibre operator sends the class of an integrable real function $f$ to
the class represented by
\[
 (I_\alpha f)(\omega)=f(\omega)\quad
 (\omega\in\Omega_h\amalg\mathcal P),\qquad
 (I_\alpha f)(\infty)=\int_T f(t)\,dt.
\]
This representative is integrable for every integrable $f$.
\end{example}
\begin{proof}
The parameter map is measurable and preserves the measure by the
probability-collapse lemma, since $T$ has measure one. On a horizontal
place every nonzero rational constant has valuation zero, so restriction
gives the trivial absolute value. On a prime place the Gauss absolute value
of a constant is its $p$-adic absolute value. On an archimedean place
evaluation fixes rational constants, so restriction gives the usual real
absolute value. At zero all of these equalities hold because absolute
values vanish at zero. This proves the absolute-value compatibility at
every parameter and every rational number. The same probability-collapse
lemma proves that the displayed function-level operator descends to a
real-linear map on integrable equivalence classes, preserves their
integrals, and retracts pullback. These facts verify every axiom of an
adelic-curve morphism.
\end{proof}
